\section{Hermitian reversal and matrix aggregation}
\label{sec:Hermitian-aggregation}

\subsection{Reversed content profiles}

Return to the actual row-pair notation.  TPC--74 proves
\[
 L_{nm,q}(v,u)
 =
 \overline{L_{mn,q}(u,v)}
 \qquad(m\ne n).
\]
The divisor set of
\[
 |\Delta_{mn}|=|h_0(n-m)|
\]
is unchanged by reversal.

\begin{proposition}[Filtered Hermitian reversal]
\label{prop:filtered-reversal}
If the reverse filter obeys
\[
 w^{nm}_{d,q}
 =
 \overline{w^{mn}_{d,q}},
\]
then
\[
 F^{W^{nm}}_{nm}(v,u)
 =
 \overline{F^{W^{mn}}_{mn}(u,v)}.
\]
In particular, the zero and zeta lifts are Hermitian compatible, and any
single real filter can be used in both orientations.
\end{proposition}

\begin{proof}
Conjugate \eqref{eq:filtered-completion}, exchange \(u,v\), and use the
literal slice reversal.
\end{proof}

Fix transpose rectangles and equal cutoffs:
\[
 B_{nm}=B_{mn}^{\mathsf T},
 \qquad
 D_{nm}=D_{mn}.
\]
The divisor--variation cost is invariant under transpose conjugation.
Thus the pair-dependent provenance gauges satisfy
\[
 \gamma^{\rm prov}_{nm}
 =
 \gamma^{\rm prov}_{mn}.
\]

\subsection{The provenance matrix}

Let \(Q^{\rm all}\) be the matrix of unrestricted TPC--75 edge gauges and
let
\[
 Q^{\rm prov}_{mn}
 =
 \begin{cases}
  \gamma^{\rm prov}_{mn},&m\ne n,\\
  0,&m=n.
 \end{cases}
\]
Let \(P^0\) and \(P^{\rm Z}\) contain the costs of the zero and zeta
completions, respectively, with zero diagonal.  More generally, for a
family of filters \(\boldsymbol W=(W^{mn})_{m<n}\), put
\[
 P(\boldsymbol W)_{mn}
 =
 \cC_{D_{mn}}\!\left(F^{W^{mn}}_{mn}\right)
 \quad(m\ne n),
 \qquad
 P(\boldsymbol W)_{mm}=0,
\]
and mirror the upper triangle.
The superscript ``all'' means all post-\(q\)-aggregation analytic
completions of the fixed primitive kernel; it does not mean a
pre-\(q\), slice-by-slice optimization.

\begin{theorem}[Four-level completion hierarchy]
\label{thm:matrix-hierarchy}
For the gauged direct-entry matrix \(\widetilde H\),
\begin{equation}
 \boxed{
 \|H\|=\|\widetilde H\|
 \le
 \rho(|\widetilde H|)
 \le
 \rho(Q^{\rm all})
 \le
 \rho(Q^{\rm prov})
 \le
 \min\{\rho(P^0),\rho(P^{\rm Z})\}.}
 \label{eq:matrix-hierarchy}
\end{equation}
\end{theorem}

\begin{proof}
Apply \cref{eq:completion-hierarchy} to every unordered row pair.
All matrices are symmetric and nonnegative, so entrywise order passes to
Perron spectral radius.  The first two inequalities are the TPC--75
absolute-entry floor.
\end{proof}

\begin{theorem}[Independent provenance factorization]
\label{thm:provenance-factorization}
If every unordered row pair may choose its own provenance filter before
the later matrix extremizer is selected, then
\[
 \inf_{\boldsymbol W}
 \rho(P(\boldsymbol W))
 =
 \rho(Q^{\rm prov}),
\]
and the infimum is attained.
\end{theorem}

\begin{proof}
Each edge minimum is attained by
\cref{thm:content-flow-dual}.  Choose all finitely many upper-triangle
minimizers simultaneously and conjugate them on the reverse edges.
Every other filter family gives an entrywise larger nonnegative
majorant, so Perron monotonicity proves equality.
\end{proof}

\begin{remark}[A universal filter is a different problem]
A single formula \(w_{d,q}\) shared across row pairs, blocks, scales, or a
finite coefficient family couples the edge variables and can be more
expensive than \(\rho(Q^{\rm prov})\).  It must be optimized with the
fixed-data affine coupled-rule theorem of TPC--75.  Uniformity over a
growing coefficient class adds a minimax quantifier and is not obtained
by choosing the best filter after each literal kernel is known.
\end{remark}

\subsection{Dual obstruction matrices}

For every edge, a feasible charge in
\cref{eq:content-flow-dual} gives
\[
 O^{\rm prov}_{mn}
 =
 \left|
 \Real\langle z_{mn},A_{D_{mn}}F^0_{mn}\rangle
 \right|
 \le
 \gamma^{\rm prov}_{mn}.
\]
Indeed, dual feasibility is invariant under \(z\mapsto-z\), so the
absolute dual objective is still bounded by the optimum.  Mirror this
value and put zero on the diagonal.

\begin{corollary}[Perron obstruction]
\label{cor:provenance-Perron-obstruction}
\[
\rho(O^{\rm prov})
\le
\rho(Q^{\rm prov}).
\]
Thus a certified nonnegligible obstruction matrix closes the declared
provenance-filter majorant at that scale.  It does not lower-bound
\(\|H\|\), because the latter may exploit signed cancellation below the
absolute Perron floor.
\end{corollary}

\begin{proof}
Use entrywise monotonicity.
\end{proof}
